\chapter{Vomiting and Nausea}

\section*{Definition}
Nausea is the unpleasant sensation of the urge to vomit, while vomiting (emesis) is the forceful expulsion of gastric contents through the mouth. Both are common symptoms with many causes, ranging from benign to life-threatening.

\section*{What Happens in Your Body}
The vomiting reflex is coordinated by the vomiting center in the medulla oblongata, which receives input from four sources: the chemoreceptor trigger zone (CTZ) in the area postrema (sensitive to toxins and drugs), the vestibular system (motion sickness), the gastrointestinal tract (via vagal afferents), and higher cortical centers (emotions, smells, pain). Activation of the vomiting center triggers a coordinated sequence: salivation, palatal elevation, diaphragmatic contraction, and gastric relaxation.

\section*{Causes}
\begin{itemize}
\item Gastrointestinal: gastroenteritis, gastritis, peptic ulcer, pancreatitis, cholecystitis, appendicitis, bowel blockage
\item Pregnancy: morning sickness (nausea gravidarum), hyperemesis gravidarum
\item Medications: chemotherapy, antibiotics, opioids, NSAIDs, anticonvulsants
\item Motion sickness
\item Inner ear disorders: labyrinthitis, Meniere disease, vestibular neuritis
\item Migraine headache
\item Metabolic: uremia, diabetic ketoacidosis, adrenal insufficiency, hypercalcemia
\item Infections: viral, bacterial, parasitic (not just GI)
\item Central nervous system: increased intracranial pressure, meningitis, tumor, concussion
\item Psychiatric: anxiety disorders, bulimia nervosa
\item Postoperative nausea and vomiting
\end{itemize}

\section*{When to See a Doctor}
\begin{itemize}
\item Nausea with or without vomiting
\item Duration and frequency of episodes
\item Vomiting blood (vomiting blood) - emergency
\item Bilious (green) or fecal-colored vomitus - suggests blockage
\item Abdominal pain, fever, or diarrhea
\item Headache, stiff neck, or photophobia
\item History of head trauma
\item Inability to keep fluids down (dehydration risk)
\item Weight loss
\end{itemize}

\section*{Related Symptoms}
\begin{itemize}
\item Abdominal pain
\item Diarrhea
\item Dehydration
\item Dizziness (with vestibular causes)
\item Headache (migraine)
\item Weight loss
\end{itemize}

\section*{Self-Care / Home Management}
\begin{itemize}
\item Eat small, frequent meals of bland foods (BRAT diet: bananas, rice, applesauce, toast)
\item Avoid strong odors, spicy or greasy foods
\item Sip clear fluids (water, clear broth, electrolyte solutions) frequently
\item Ginger (tea, candied, or supplements) can reduce nausea
\item Peppermint aromatherapy
\item Acupressure wristbands (P6 point)
\item Rest and avoid sudden head movements
\end{itemize}

\section*{How Your Doctor Will Evaluate You}
\begin{itemize}
\item History and physical examination
\item Pregnancy test in women of reproductive age
\item Basic metabolic panel, liver enzymes, amylase, lipase
\item Abdominal imaging (X-ray, CT, ultrasound) based on suspected cause
\item Upper endoscopy if upper GI source suspected
\end{itemize}

\section*{Treatment Options}
\begin{itemize}
\item Antiemetic medications: ondansetron, promethazine, metoclopramide, prochlorperazine
\item IV fluids for dehydration
\item Treat underlying cause (antibiotics for infection, migraine therapy, chemotherapy antiemetics)
\item Electrolyte replacement
\item Nutritional support for prolonged symptoms
\end{itemize}
